\section{Probability Revision}

This section revisits the formulation of the probability of the emergence of life. The aim is not to assign a numerical value, but to clarify what the probability is taken to measure.

\subsection{Conventional Formulation}

The emergence of life is often treated as the probability that a particular molecular structure appears by chance. In this view, many conditions must be satisfied simultaneously, and each may have low probability. The combined probability can therefore become very small.

This view is coherent if life is assumed to be a particular completed configuration. The present framework uses a different object of probability.

\subsection{The Issue}

As described in Section 7, life is not a fixed structure. It is sustained stable motion in state space. Life therefore corresponds not to a single point, but to a region with finite extent.

The emergence of life is then interpreted as reaching this region and remaining within it, rather than hitting one specified structure.

\subsection{Probability as Region Entry}

Let the life attractor occupy a region
\begin{equation}
  \Omega_{\mathrm{life}} \subset \Omega.
\end{equation}
The probability of life emergence is then the probability of reaching \(\Omega_{\mathrm{life}}\). This is a probability of region entry, not a probability of exact point matching.

\subsection{Structure of Trials}

Molecular reactions and material cycles repeat as long as enabling conditions persist. States that do not sustain themselves disappear. States that sustain themselves remain. In this sense, life appears not because it is selected in advance, but because certain states can persist.

\subsection{Dimensional Reduction}

In a point-target view, many degrees of freedom must be specified simultaneously. In the present framework, the relevant conditions for life are described by fewer effective variables that characterize entry into a stable region. The dimensional structure of the probability problem is therefore changed.

\subsection{Interpretation}

This is not a numerical reassessment of probability. It is a structural redefinition of the object being counted. Life is neither a miraculous hit nor an automatic necessity. It is a state that can appear when sustaining conditions exist.

In IFGT terms, this reformulation has a precise interpretation: life corresponds to the entry and persistence of a system within a quasi-closure regime, a dynamical region in which the drive toward closure is sustained while complete closure remains unattained as an ordinary dynamical state. The probability of life is therefore better understood as the probability of reaching and maintaining such a regime, rather than the probability of assembling a fixed molecular configuration. The apparent improbability of life may reflect, in part, a consequence of how the target was identified.

\subsection{Connection to Intelligence}

This reformulation makes the next question more precise: after life is established, why do only some systems reach intelligence? Intelligence is not a simple continuation of life. It requires additional structural conditions. The next section treats intelligence as a second attractor.
